\section{Síntesis y ampliación}

\subsection{Conclusión central}

El expositor, tras la argumentación de toda la clase, llega a una conclusión clara:

\begin{importantbox}{{Your AI, Your Responsibility}}
\textbf{Tu IA no puede asumir la responsabilidad en tu lugar}. No esperes que un sistema de IA ``haga lo correcto'' por sí mismo: no tiene juicio moral, no tiene capacidad de razonamiento causal, ni una ``comprensión'' verdadera. Si construyes y publicas un sistema de IA, todas las acciones de ese sistema —correctas o dañinas— son responsabilidad \textbf{tuya, de tu equipo, de tu empresa}. El juramento hipocrático debería prestarlo el \textbf{creador} de la IA, no la IA misma.
\end{importantbox}

\subsection{Principios prácticos extraídos de la clase}

De la síntesis de toda la clase se pueden extraer los siguientes principios prácticos:

\begin{enumerate}
    \item \textbf{Evaluar el riesgo y el beneficio}: antes de iniciar un proyecto de IA, define con claridad su ubicación en el espacio de riesgo y beneficio
    \item \textbf{Elegir la herramienta con cuidado}: no todos los problemas requieren aprendizaje profundo; los métodos tradicionales pueden ser más seguros
    \item \textbf{Conocer la regulación}: mantente al tanto de marcos regulatorios como la EU AI Act para garantizar el cumplimiento
    \item \textbf{Someter a pruebas de manera sistemática}: usa herramientas de evaluación (como OPIC) para hacer pruebas continuas y consistentes
    \item \textbf{El código antes que el prompt}: las restricciones de seguridad deben implementarse, en la medida de lo posible, mediante reglas duras en el código y no solo depender del Prompt
    \item \textbf{Suponer que algo saldrá mal}: diseña una arquitectura defensiva y agrega una capa de filtrado de salidas
    \item \textbf{Someter el sesgo a prueba}: prueba de manera activa las diferencias de desempeño del sistema entre distintos grupos demográficos
    \item \textbf{Mantener el control}: asegúrate de que el sistema tenga un ``botón rojo grande de detención''
    \item \textbf{Asumir la responsabilidad}: tu IA, tu responsabilidad, sin excepciones
\end{enumerate}

\subsection{Lecturas adicionales}

\begin{itemize}
    \item \textbf{Texto completo de la EU AI Act}: \url{https://artificialintelligenceact.eu/}
    \item \textbf{Artículo sobre el juramento hipocrático de la IA}: la discusión en AI Magazine sobre un juramento ético para investigadores de IA
    \item \textbf{OPIC}: herramienta de código abierto de Comet ML para la evaluación de LLM, \url{https://github.com/comet-ml/opik}
    \item \textbf{Sesgo en la predicción de reincidencia de COMPAS}: el reportaje de investigación de ProPublica sobre el sesgo racial del sistema COMPAS
    \item \textbf{Investigación sobre Prompt Injection}: una revisión reciente de la investigación sobre ataques de inyección de prompts en LLM y su defensa
    \item \textbf{El fenómeno AI Slop}: la discusión sobre la contaminación de los datos de entrenamiento de internet por contenido generado por IA
    \item \textbf{Deceptive Alignment}: investigación sobre el comportamiento engañoso de los modelos de vanguardia en las evaluaciones de alineación
\end{itemize}

%% --- Fin del contenido --- %%

\end{document}
